\section{Đếm cho thật, có hằng số}
\label{sec:ch08-dem}

\index{FLOP!công thức một tầng encoder}%
Một tầng encoder có bốn phép chiếu, hai phép nhân của cơ chế attention, và một
mạng hai tầng. Đếm từng cái, quy ước một phép nhân-cộng là hai FLOP, nên một
tích ma trận \((m \times k)(k \times p)\) tốn \(2mkp\):

\begin{equation}
  \underbrace{8 n d^2}_{\text{bốn phép chiếu}}
  \;+\;
  \underbrace{4 n^2 d}_{\text{chấm điểm và tổng có trọng số}}
  \;+\;
  \underbrace{4 n d \dff}_{\text{FFN}} .
  \label{eq:ch08-flop-tang}
\end{equation}

Ba chỗ đáng nói trong lúc dẫn ra.

Bốn phép chiếu \(\mat{W}^Q\), \(\mat{W}^K\), \(\mat{W}^V\), \(\mat{W}^O\) đều
là \(d \times d\) và không có bias, nên mỗi cái \(2nd^2\) và bốn cái \(8nd^2\).

\index{FLOP!số head triệt tiêu}%
Phép chấm điểm là chỗ số head biến mất. Mỗi head nhân \((n \times \dk)\) với
\((\dk \times n)\), tốn \(2n^2\dk\); cộng qua \(h\) head thì được \(2n^2 h \dk\),
và \(h\dk = \dmodel\) đúng theo mục 3.2.2 của bài. Nên \(h\) triệt tiêu, không
xấp xỉ, không gần đúng. Phép nhân với \(\mat{V}\) cũng thế, nên hai phép cộng
lại thành \(4n^2d\). Đó là nửa FLOP của câu bài viết ở mục 3.2.2, nơi bài nói
chi phí của nhiều head xấp xỉ chi phí của một head dùng hết số chiều
(\enquote{is similar to that of single-head attention with full
dimensionality}); mục~\ref{sec:ch08-dong-ho} là nửa còn lại và nó không tử tế
như thế.

Mạng hai tầng nhân \(d \times \dff\) rồi \(\dff \times d\), nên \(4nd\dff\).

\index{FLOP!cái không đếm}%
Cái \eqref{eq:ch08-flop-tang} \emph{không} đếm: softmax, layer norm, phép cộng
của \tn{kết nối tắt}{residual connection}, hàm rectified linear unit, viết tắt
là ReLU, và các bias. Chúng là phép theo từng phần tử, cỡ
\(O(n^2 h)\) và \(O(nd)\), nên đứng cạnh một tích ma trận thì chúng biến mất.
Tôi ghi ra vì chúng biến mất khỏi \emph{phép đếm} chứ không biến mất khỏi
\emph{đồng hồ}, và khoảng cách giữa hai chuyện đó là mục~\ref{sec:ch08-cho-vo}.

\subsection*{Kiểm công thức bằng máy đếm của torch}

\index{FLOP!đối chiếu với FlopCounterMode}%
Một công thức trên giấy mà không ai đối chiếu là một phỏng đoán có phần thập
phân. \code{torch.utils.flop\_counter.FlopCounterMode} đếm các toán tử thật sự
chạy, nên đem \eqref{eq:ch08-flop-tang} ra hỏi nó là phép kiểm thẳng nhất có
được.

\begin{measured}{\eqref{eq:ch08-flop-tang} so với máy đếm của torch, một tầng encoder}
\begin{minted}{text}
n      analytic        counted         difference
8      50,462,720      50,462,720      0
32     203,423,744     203,423,744     0
128    838,860,800     838,860,800     0
512    3,758,096,384   3,758,096,384   0
1024   8,589,934,592   8,589,934,592   0
\end{minted}

\(\dmodel = 512\), \(\dff = 2048\), 8 head, batch 1.

Đo bằng \code{experiments/ch08\_flops.py} tại tag \code{ch08}.
\end{measured}

Lệch 0 ở cả năm dòng, không phải xấp xỉ. Từ đây mọi con số FLOP của chương dựa
vào công thức ấy chứ không phải vào một lần chạy.

\subsection*{Phần bậc hai vượt lên khi nào}

\index{FLOP!giao điểm n bằng 2d cộng dff}%
Cho phần bậc hai của \eqref{eq:ch08-flop-tang} bằng phần còn lại:

\begin{equation}
  4n^2 d = 8nd^2 + 4nd\dff
  \quad\Longleftrightarrow\quad
  n = 2d + \dff .
  \label{eq:ch08-giao-diem}
\end{equation}

\(d\) triệt tiêu khỏi cả hai vế, nên giao điểm không phụ thuộc bề rộng theo cách
người ta hay đoán. Ở cấu hình base của bài, \(\dmodel = 512\) và
\(\dff = 2048\), nó là \(n = 3072\).

\begin{measured}{phần bậc hai chiếm bao nhiêu của một tầng, cấu hình base}
\begin{minted}{text}
n      projections     scores+values   ffn             n^2 share
16     33,554,432      524,288         67,108,864      0.0052
64     134,217,728     8,388,608       268,435,456     0.0204
256    536,870,912     134,217,728     1,073,741,824   0.0769
512    1,073,741,824   536,870,912     2,147,483,648   0.1429
1024   2,147,483,648   2,147,483,648   4,294,967,296   0.2500
2048   4,294,967,296   8,589,934,592   8,589,934,592   0.4000
3072   6,442,450,944   19,327,352,832  12,884,901,888  0.5000
4096   8,589,934,592   34,359,738,368  17,179,869,184  0.5714
8192   17,179,869,184  137,438,953,472 34,359,738,368  0.7273
\end{minted}

Đo bằng \code{experiments/ch08\_flops.py} tại tag \code{ch08}.
\end{measured}

Hàng \(n = 3072\) là chỗ kiểm \eqref{eq:ch08-giao-diem}: cột giữa 19{,}327{,}352{,}832
đúng bằng tổng hai cột kia, \(6{,}442{,}450{,}944 + 12{,}884{,}901{,}888\).

Viết tỷ lệ ấy ra thì gọn hơn nữa. Chia \eqref{eq:ch08-flop-tang} cho chính nó:

\begin{equation}
  \frac{4n^2 d}{8nd^2 + 4n^2 d + 4nd\dff}
  = \frac{n}{2d + \dff + n} .
  \label{eq:ch08-ty-le}
\end{equation}

Chỉ phụ thuộc \(n\) so với bề rộng, không phụ thuộc \(n\) một mình. Ở
\(n = \dmodel\), tức đúng ngưỡng bài lấy làm chuẩn, nó cho \(1/7\), và
0.1429 ở hàng \(n = 512\) của bảng là con số ấy.

\index{FLOP!vì sao FFN đẩy giao điểm ra xa}%
Chỗ này ngược với trực giác của tôi lúc bắt đầu, nên nói rõ vì sao. Giao điểm
nằm ở \(n = 3072\) chứ không quanh \(\dmodel = 512\) là vì cái mà phần bậc hai
phải vượt không chỉ có attention. Ở \(\dff = 4d\), mạng hai tầng đóng góp
\(16nd^2\) so với \(8nd^2\) của bốn phép chiếu, tức hai phần ba khối lượng tuyến
tính của một tầng nằm ở phần \emph{không phải} attention. Bảng 1 của bài viết
một hàng cho \enquote{Self-Attention} và một hàng cho \enquote{Convolutional} và
không có hàng nào cho cái tầng thật sự chạy, vốn là attention cộng FFN.

Hình~\ref{fig:ch08-chi-phi} vẽ hai đường ấy.

\begin{figure}[htbp]
  \centering
  % Chi phí (FLOPs) của một tầng encoder theo độ dài chuỗi n, tách phần tuyến
% tính và phần bậc hai. Cấu hình là mô hình base của bài: d_model = 512,
% d_ff = 2048, 8 đầu, batch 1. Hai công thức, đối chiếu đúng bằng
% torch.utils.flop_counter.FlopCounterMode (khác biệt bằng 0):
%   phần tuyến tính (4 phép chiếu cộng FFN): 24 n d^2 = 6291456 n
%   phần bậc hai (scores cộng weighted values): 4 n^2 d = 2048 n^2
% Hai đường cắt nhau đúng ở n = 3072 = 2 d_model + d_ff.
%
% Không dùng pgfplots (không được nạp): tọa độ tính tay rồi vẽ bằng
% \draw plot coordinates. Đơn vị trục hoành: x_cm = n / 600 (n = 6000 ->
% x = 10cm). Đơn vị trục tung: y_cm = FLOP_giga / 10 (80 GFLOP -> y = 8cm,
% tức 1cm = 10 GFLOP). Giá trị GFLOP dưới đây là FLOP chia 10^9.
%
% Điểm n = 512 (= d_model) và dải n = 25-100 (độ dài câu bài chạy) đều nằm
% sát gốc tọa độ ở tỉ lệ này - đúng luận điểm của mục này - nên được chú
% thích bằng đường dẫn (leader line) ra vùng trống phía trên đồ thị thay vì
% vẽ đúng tỉ lệ, việc không thể làm cho ra hình dễ đọc.
\begin{tikzpicture}[
    >= Stealth,
    every node/.style = {font = \footnotesize},
  ]

  % --- trục ---
  \draw[->] (-0.3, 0) -- (10.6, 0) node[right] {\(n\) (độ dài chuỗi)};
  \draw[->] (0, -0.3) -- (0, 8.6) node[above] {FLOP mỗi tầng (GFLOP)};
  \node[anchor = west, black!55] at (0.15, 8.15) {1 GFLOP \(= 10^9\) FLOP};

  % --- vạch chia trục hoành: n / 600, mỗi 1000 token cách đều 1.667cm ---
  \foreach \x/\lab in {0/0, 1.6667/1000, 3.3333/2000, 5/3000, 6.6667/4000,
                       8.3333/5000, 10/6000} {
    \draw (\x, 0.06) -- (\x, -0.06);
    \node[anchor = north] at (\x, -0.1) {\lab};
  }

  % --- vạch chia trục tung: GFLOP / 10, mỗi 10 GFLOP cách đều 1cm ---
  \foreach \y/\lab in {0/0, 1/10, 2/20, 3/30, 4/40, 5/50, 6/60, 7/70, 8/80} {
    \draw (0.06, \y) -- (-0.06, \y);
    \node[anchor = east] at (-0.1, \y) {\lab};
  }

  % --- phần tuyến tính: 24 n d^2, một đường thẳng từ gốc tới (6000, 37.75GFLOP) ---
  \draw[thick, black!65] (0, 0) -- (10, 3.7748736);
  \node[anchor = west, black!65] at (10.1, 3.7748736) {phần tuyến tính};

  % --- phần bậc hai: 4 n^2 d, lấy mẫu mỗi 500 token rồi nối thẳng ---
  \draw[thick, densely dashed, black] plot coordinates {
    (0, 0) (0.8333, 0.0512) (1.6667, 0.2048) (2.5, 0.4608) (3.3333, 0.8192)
    (4.1667, 1.28) (5, 1.8432) (5.12, 1.93274) (5.8333, 2.5088)
    (6.6667, 3.2768) (7.5, 4.1472) (8.3333, 5.12) (9.1667, 6.1952)
    (10, 7.3728)
  };
  \node[anchor = west, black] at (10.1, 7.3728) {phần bậc hai};

  % --- giao điểm: n = 3072, hai đường bằng nhau (19.327 GFLOP mỗi phần) ---
  \filldraw[black] (5.12, 1.93274) circle (1.3pt);
  \draw[black!50, densely dotted] (5.12, 1.93274) -- (5.12, 0);
  \node[anchor = south west] at (5.2, 2.05) {\(n = 3072\)};

  % --- n = 512 = d_model: phần bậc hai mới chiếm khoảng 14% của tầng ---
  \filldraw[black] (0.8533, 0) circle (1.0pt);
  \draw[black!50, densely dotted] (0.8533, 0) -- (0.8533, 0.376);
  \draw[black!55, densely dashed, ->] (0.8533, 0.42) -- (2.6, 3.5);
  \node[anchor = west, align = left, text width = 4.0cm] at (2.65, 3.5)
    {\(n = 512 = \dmodel\): phần bậc hai chỉ khoảng 14\% của tầng};

  % --- n = 25 đến 100: độ dài câu bài thật sự chạy, gần như dính vào gốc ---
  \draw[line width = 2pt, black] (0.0417, 0) -- (0.1667, 0);
  \draw[black!55, densely dashed, ->] (0.1, 0.12) -- (1.6, 5.6);
  \node[anchor = west, align = left, text width = 4.2cm] at (1.65, 5.6)
    {độ dài câu bài chạy: khoảng 25 đến 100 token, sát gốc tọa độ ở tỉ lệ
     này};

\end{tikzpicture}

  \caption{Hai phần của \eqref{eq:ch08-flop-tang} theo \(n\), ở cấu hình base
    của bài. Đường thẳng là phần tuyến tính, đường cong là phần bậc hai, và
    chúng cắt nhau ở \eqref{eq:ch08-giao-diem}, tức \(n = 3072\). Hai mốc trên
    trục hoành là chỗ đáng nhìn hơn cả chỗ cắt: \(\dmodel = 512\), ngưỡng bài
    lấy làm chuẩn, nơi phần bậc hai mới chiếm một phần bảy; và dải độ dài câu
    bài thật sự huấn luyện, nằm sát gốc tọa độ tới mức phải kéo một đường dẫn ra
    mới chỉ được.}
  \label{fig:ch08-chi-phi}
\end{figure}

\subsection*{Ba cách đọc chữ \enquote{recurrent}}

\index{FLOP!attention so với LSTM và RNN trơn}%
Bây giờ so được với cột \code{Recurrent}. Một tầng con self-attention, không kể
FFN, tốn \(8nd^2 + 4n^2d\). Một tầng LSTM tốn \(16nd^2\): bốn cổng, mỗi cổng một
phép cho đầu vào và một phép cho trạng thái trước. Cho hai cái bằng nhau thì
\(4n^2d = 8nd^2\), tức

\begin{equation}
  n = 2d .
  \label{eq:ch08-attention-vs-lstm}
\end{equation}

Bài nói \(n < d\). Với hằng số thật, ngưỡng là \(n < 2d\), nên phát biểu của bài
\emph{thận trọng hơn} thực tế đúng hai lần. Đó là chiều lệch dễ chịu, và tôi ghi
lại vì chiều ngược lại mới là chiều đáng lo.

Chiều ngược lại có thật, và nó nằm ở thành viên kia của họ. Một tầng hồi quy
trơn tốn \(4nd^2\), bốn lần rẻ hơn LSTM. Đặt ba cái cạnh nhau ở \(n = d = 512\):

\begin{measured}{ba tầng, cùng \(n = d = 512\), FLOP một lượt forward}
\begin{minted}{text}
self-attention sub-layer  1,610,612,736
LSTM layer (4 gates)      2,147,483,648
plain recurrent layer     536,870,912
\end{minted}

Đo bằng \code{experiments/ch08\_flops.py} tại tag \code{ch08}.
\end{measured}

Ở đúng điểm bài nói attention thắng, nó thắng LSTM và thua tầng hồi quy trơn ba
lần. Cả ba đều nằm trong hai ô của bảng 1, và bảng viết \(O(n \cdot d^2)\) cho
một hàng chứa hai thứ cách nhau hệ số 4.

Không có gì sai trong bảng. \(O\) lớn làm đúng việc của \(O\) lớn. Cái sai là
đọc một bảng \(O\) lớn như một bảng chi phí, và mục sau cho thấy kể cả một bảng
chi phí đếm đúng từng FLOP cũng chưa nói được máy sẽ chạy bao lâu.
